\documentclass[11pt,a4paper]{article}

\usepackage[T1]{fontenc}
\usepackage{lmodern}
\usepackage[a4paper,margin=25mm]{geometry}
\usepackage{amsmath,amssymb,amsthm,mathtools}
\usepackage{booktabs,array}
\usepackage{microtype}
\usepackage{xcolor}
\usepackage{enumitem}
\usepackage{placeins}
\usepackage{url}
\usepackage[colorlinks=true,linkcolor=blue!55!black,citecolor=blue!55!black,
            urlcolor=blue!55!black]{hyperref}

\newtheorem{theorem}{Theorem}
\newtheorem{lemma}[theorem]{Lemma}
\newtheorem{proposition}[theorem]{Proposition}
\newtheorem{remark}[theorem]{Remark}
\newcommand{\Z}{\mathbb{Z}}
\newcommand{\R}{\mathcal{R}}
\newcommand{\E}{\mathbb{E}}
\newcommand{\Prb}{\mathop{\rm Pr}}
\newcommand{\Sign}{\mathsf{Sign}}
\newcommand{\Verify}{\mathsf{Verify}}
\newcommand{\pk}{\mathsf{pk}}
\newcommand{\sk}{\mathsf{sk}}
\newcommand{\norm}[1]{\left\lVert #1\right\rVert}

\newcommand{\RealRows}{1{,}934{,}304}
\newcommand{\RealCollectionSeconds}{5{,}207.6}
\newcommand{\RealCollectionRate}{576.1}
\newcommand{\RealRadialRMSE}{0.244287}
\newcommand{\RealRadialMaxError}{0.884067}
\newcommand{\RealRoundedExact}{731}
\newcommand{\RealNeighborSurvivors}{0}
\newcommand{\RealRadialSeconds}{86.2}
\newcommand{\RealIntegerSeconds}{125.1}
\newcommand{\RealLatticeSeconds}{480.3}
\newcommand{\RealPeakRSS}{1.45~GiB}
\newcommand{\RealLWRRange}{\ensuremath{[-16,15]}}
\newcommand{\RealRecoveredExact}{768}
\newcommand{\RecoveryHash}{aacacca14b8fb842670ea85e939adc0229f04a212aec34d57af50ba7150c4669}
\newcommand{\ForgeryHash}{cda3afd63f3f1149b2b1ffe7746ae4db045dbf77bdbf88b1e92eda5ebd838c29}
\newcommand{\EquivalentKeyHash}{fd3bdc055d0bb376031bf556721c371c7e0c2b6f56850fc6aaf6b3509e2e499b}
\newcommand{\ProspectiveRows}{1{,}931{,}974}
\newcommand{\ProspectiveCollectionSeconds}{5{,}064.7}
\newcommand{\ProspectiveCollectionRate}{592.3}
\newcommand{\ProspectiveRadialRMSE}{0.249334}
\newcommand{\ProspectiveRadialMaxError}{0.833201}
\newcommand{\ProspectiveRoundedExact}{728}
\newcommand{\ProspectiveNeighborSurvivors}{0}
\newcommand{\ProspectiveRadialSeconds}{160.0}
\newcommand{\ProspectiveIntegerSeconds}{126.1}
\newcommand{\ProspectiveLatticeSeconds}{75.6}
\newcommand{\ProspectivePeakRSS}{1.42~GiB}
\newcommand{\ProspectiveRecoveryHash}{c5111d7b13af090288d583081f1f493747bd28283042069e9d40ccffe8c73375}
\newcommand{\ProspectiveForgeryHash}{b4a98b0913b409c8ded21762cf0bd69b4d08aef19fba4e7f2c4c85a6869da2bb}
\newcommand{\ProspectiveEquivalentKeyHash}{32130947f97b0982c4e89124b7f0f076f926042810f9e35bc579d610469f49d9}
\newcommand{\ProspectiveTranscriptHash}{355829fab9c9b869fa0d28eebd0e93258d276f5be69149228aeeecbaf9e7bec4}
\newcommand{\FreezeManifestHash}{9bebb94f64c14906f1284b25293a32d65bb86504c8bc96f95300330b0545a18c}
\newcommand{\PublicOutcomeHash}{f89c9e1433ea6ecca1ae963b40bee4992e5eaa09c56587c2419983751973ecfd}

\title{\bfseries First-Gate Support Leakage in MORNING-ATLAS-128:\\
Exact $s_1$ Recovery and Fresh-Message Forgeries\\
from Three Million Signatures per Key}
\author{Martin Feussner\\
\small Selmer Center, University of Bergen\\
\small \texttt{martin.feussner@uib.no}}
\date{27 September 2026}

\begin{document}
\maketitle

\begin{center}
\fbox{\parbox{0.92\linewidth}{\small\textbf{Provenance and review status.}
The attack was independently found by an AI system, OpenAI Codex using
Daybreak Blue at maximum reasoning effort, without a human first proposing
this specific attack.  The derivation, implementation, experiments, and
initial manuscript were also produced by the AI.  At the date above, no human
cryptanalyst has independently verified the argument, experiments, or novelty
assessment.  This is a preliminary disclosure for human review and
reproduction.}}
\end{center}
\vspace{0.7em}

\begin{abstract}
MORNING-ATLAS is a module-LWR signature submitted to the Next-generation
Commercial Cryptographic Algorithms Program.  ATLAS-128 samples every mask
coefficient uniformly from $[-(\gamma_1-1),\gamma_1-1]$ and publishes
$z=y+cs_1$ only when $\norm{z}_\infty<\gamma_1-\beta_1$.  The specification
argues that translating the accepted interval by a coefficient of $cs_1$
does not change its size.  This overlooks truncation by the finite support of
$y$.  For ATLAS-128, $\beta_1=51$ while a challenge product can have magnitude
$\kappa\eta=496$.  If $a=(cs_1)_j$, the number of masks compatible with an
accepted response coefficient is exactly
\[
  1{,}048{,}473-(|a|-51)_+,
\]
and boundary-near response coefficients supply exact public half-spaces containing $s_1$.

Feasibility alone is insufficient because the zero vector satisfies all of
these half-spaces.  In an ideal first-gate model, however, the conditional
likelihood of every feasible transcript is inversely proportional to the
candidate's first-gate acceptance probability.  A public bound-moment
regression and radial linear programs put every target coefficient inside a
public floor/ceiling box.  Response-only integer programs give a feasible
baseline.  A deterministic polynomial-support ladder then combines 64 public
LWR equations with each binary box in a Kannan embedding; LLL followed by
BKZ-10 exposes $s_1$ exactly.  A public-hint step recovers the
omitted low part $t_0$, and the recovered $(s_1,t_0)$ forms an equivalent
signing key.

On a calibration key, an execution on 3,000,000 ordinary chosen-message
signatures retained \RealRows{} support rows, recovered all 768 coefficients,
and produced a verifier-accepted fresh-message forgery.  Because the
completion method was developed retrospectively on that key, we then froze
the code, parameters, public block order, and stopping rules before selecting
a second key.  With both ground-truth secret components withheld from the recovery and
evaluation pipeline until the public result and forgery were sealed, the same
procedure recovered exact $(s_1,t_0)$ from a new three-million-signature
transcript and forged a fresh message.  On this prospective run, the response-only integer baseline was
already exact; the prescribed lattice stage rediscovered it at block 0 after
BKZ-10.  Four-core collection took \RealCollectionSeconds{} and
\ProspectiveCollectionSeconds{} seconds on the two keys, and peak public
recovery memory was \RealPeakRSS{}.  The corrected samplers target the
written design, separately from the published two-signature repeated-mask
implementation attack.  Both tested keys succeeded, but two trials do not
estimate a key-averaged success probability.
\end{abstract}

\section{Introduction}

Fiat--Shamir-with-aborts signatures publish a response obtained by adding a
secret-dependent challenge product to a random mask.  Rejection sampling is
supposed to remove the dependence on the signing secret.  A proof of
independence must account for both the desired response interval and the
finite support from which the mask was originally sampled.

MORNING-ATLAS follows this paradigm over a power-of-two cyclotomic ring
\cite{ATLAS}.  Its specification samples $y$ uniformly from a coefficient
box and forms $z=y+cs_1$.  The first signing gate retains $z$ only inside a
slightly smaller box.  Section~2.5 of the specification states that shifting
that smaller interval by a coefficient of $cs_1$ leaves exactly
$2(\gamma_1-\beta_1)-1$ admissible mask values.  That count is valid only when
the shifted interval stays inside the sampler support.  The ATLAS-128
parameters do not ensure this condition: $\beta_1=51$ and
$\norm{cs_1}_\infty$ can be as large as $496$.

This mismatch has two consequences.  First, a response near a mask boundary
gives an exact linear inequality in $s_1$.  Second, challenges that produce
larger products have lower first-gate acceptance probability.  Once we
condition on observing a valid signature, feasible candidate secrets with
lower acceptance probability receive higher likelihood.  A direct likelihood
optimization is nonlinear.  A Gaussian approximation suggests that the
acceptance loss grows with the squared coefficient norm, leading to a simple
sequence of public linear programs over the exact support polytope.

The public key supplies a decisive validation relation.  If the full rounded
product $t$ is known, a candidate $s_1$ must satisfy
\[
            A s_1-32t=e\pmod{2^{23}},
            \qquad e\in[-16,15]^{9\times128}.
\]
ATLAS omits the low part $t_0$ from its public key, but ordinary signature
hints expose bounded observations of $ct_0$.  A public minimax computation on
20,000 signatures recovers all 1,152 coefficients of $t_0$.  The same
messages are a subset of the three-million-signature response transcript, so
the total number of distinct chosen-message queries remains three million.

Our contributions are:

\begin{itemize}[leftmargin=1.5em]
  \item We give the exact accepted-support formula and identify the omitted
        truncation in the ATLAS rejection-sampling analysis.
  \item We derive exact public half-spaces and the ideal first-gate candidate
        likelihood.  We clearly separate these statements from the
        norm-based heuristic used for efficient recovery.
  \item We specify a recovery procedure without access to the target secret,
        consisting of a public regression, radial cutting-plane linear
        programs, response-feasible integer completion, and a deterministic
        LWR lattice ladder with complete public validation.
  \item On two separately generated ATLAS-128 targets, the procedure recovers
        $s_1$ exactly, combines it with publicly recovered $t_0$, constructs
        equivalent signing keys, and produces accepted signatures on fresh
        messages.  The second execution was frozen before its key was
        generated and completed before the target secret was revealed to the
        recovery and evaluation process.
  \item We test the written design with corrected uniform samplers and
        distinguish the result from the previously published implementation
        attack caused by repeated adjacent mask coefficients \cite{NGCC15}.
\end{itemize}

The attack uses no timing information, fault, malformed signature, submitted
sampler repetition, or access to signing randomness.  It consumes fields
decoded from signatures returned by an ordinary chosen-message signing
oracle.

\section{ATLAS-128 and the faithful target}
\label{sec:atlas}

Let
\[
  \R_q=\Z_q[X]/(X^{128}+1),\qquad q=2^{23},\qquad p=2^{18}.
\]
ATLAS-128 uses a public matrix $A\in\R_q^{9\times6}$ and a secret
$s_1\in\R_q^6$ whose 768 coefficients are independently uniform in
$[-\eta,\eta]$ with $\eta=16$.  Key generation rounds $As_1$ from modulus
$q$ to modulus $p$.  Writing that rounded vector as $t\in\R_p^9$ gives the
centered public relation
\begin{equation}
                 As_1-32t=e\pmod q,
        \qquad e\in[-16,15]^{9\times128}.
\label{eq:lwr}
\end{equation}
The public key stores a matrix seed and $t_1$; with $d=10$,
\begin{equation}
                    t=2^d t_1+t_0\pmod p,
\label{eq:tdecomp}
\end{equation}
where every centered coefficient of the omitted $t_0$ belongs to
$[-511,512]$.  The packed signing key contains $s_1$ and $t_0$, together with
the public matrix seed, a signing-randomness seed, and a hash of the public
key.

For each signing attempt, the signer samples
\[
       y\leftarrow[-(\gamma_1-1),\gamma_1-1]^{6\times128},
       \qquad \gamma_1=2^{19},
\]
derives a signed sparse challenge $c$ of Hamming weight $\kappa=31$, and
forms
\begin{equation}
                         z=y+cs_1.
\label{eq:response}
\end{equation}
The first gate restarts unless
\begin{equation}
                 \norm{z}_\infty<\gamma_1-\beta_1,
                 \qquad \beta_1=51.
\label{eq:firstgate}
\end{equation}
Later gates check a second low-part condition, consistency of decomposed high
bits, the magnitude of $ct_0$, and the hint weight.  The signature transmits
$z$, the hint, and $c$.

\subsection{Separating the written design from submitted-code defects}

The submitted ATLAS-128 sampler has two defects relevant to any experiment.
Its mask decoder overwrites the first 20-bit value before storing it, so every
pair of adjacent mask coefficients is equal.  Lu and Liu used that relation
to recover $s_1$ from two signatures and forge a new message; the NGCC report
classifies this as an implementation issue \cite{NGCC15}.  Separately, the
submitted $\eta=16$ secret sampler splits each byte into two four-bit values,
so it cannot sample the specified 33-element interval and in fact produces
only positive coefficients.

Our faithful experiment replaces those routines with rejection samplers for
two independent 20-bit mask values and for a uniform six-bit encoding of
$[-16,16]$.  It leaves the key equations, signing gates, packing, challenge
generation, and verifier logic unchanged.  Thus adjacent-mask cancellation
is unavailable, and the target key follows the distribution stated in the
specification.

\section{Exact first-gate support leakage}
\label{sec:support}

Consider one response coefficient and write
\[
  a=(cs_1)_j,\qquad z=y+a,\qquad
  Y=\gamma_1-1=524287,
\]
and
\[
  C=\gamma_1-\beta_1-1=524236.
\]
The mask support is $[-Y,Y]$, and the strict norm test in
Equation~\eqref{eq:firstgate} accepts precisely $z\in[-C,C]$.

\begin{lemma}[Accepted support]
\label{lem:support}
For fixed $a$, the accepted response support and its cardinality are
\begin{align}
 I(a)&=[-C,C]\cap[a-Y,a+Y],\label{eq:interval}\\
 |I(a)|&=1{,}048{,}473-(|a|-51)_+.\label{eq:width}
\end{align}
\end{lemma}

\begin{proof}
The response must satisfy both $-C\le z\le C$ and
$-Y\le z-a\le Y$, which gives Equation~\eqref{eq:interval}.  If
$|a|\le Y-C=51$, the first interval is contained in the second and has
$2C+1=1{,}048{,}473$ elements.  Translating farther removes exactly
$|a|-51$ integers from one endpoint.  The maximum possible magnitude
$\kappa\eta=496$ is far below the point at which both endpoints would be
truncated, so Equation~\eqref{eq:width} follows.
\end{proof}

The specification's constant-width calculation is therefore valid only for
$|a|\le51$.  It separately notes that the parameter-selection constraint
forces $\beta_1\approx\kappa\eta$ to cease being true, but the rejection
analysis does not propagate that fact.  In particular, the claim that the
response distribution is statistically independent of $s_1$ does not follow
for ATLAS-128.

Every observed coefficient gives
\begin{equation}
                     z-Y\le a\le z+Y.
\label{eq:rawconstraint}
\end{equation}
Since $a$ is a signed negacyclic linear form in the 128 coefficients of one
secret polynomial, Equation~\eqref{eq:rawconstraint} is public and linear.
The a priori bound $|a|\le\kappa\eta=496$ makes most instances vacuous.  A
nonvacuous lower inequality has the normalized form
\begin{equation}
                    \langle x,s\rangle\ge b,
                    \qquad b\in\{-495,\ldots,-51\}.
\label{eq:normalized}
\end{equation}
A nonvacuous upper inequality is negated into the same form.  At most one
such row arises from a response coefficient.  In the three-million ideal
calibration, approximately 0.647 rows arose per signature.

\begin{remark}[Why feasibility is insufficient]
Every bound in Equation~\eqref{eq:normalized} is negative, so the all-zero
secret satisfies every support inequality.  A feasibility solve or an
ordinary minimum-norm decoder is consequently directed away from the true
secret.
\end{remark}

\section{Conditional likelihood and a radial surrogate}
\label{sec:likelihood}

Let $G$ be the first-gate event.  In an ideal random-oracle model, a fresh
challenge is uniform over the signed weight-31 challenge set before the gate.
For a candidate $s$, define
\begin{equation}
 P_s(G)=\E_c\left[
     \prod_{r=1}^{6}\prod_{j=0}^{127}
       \frac{|I((cs_r)_j)|}{2Y+1}
                  \right].
\label{eq:acceptance}
\end{equation}

\begin{proposition}[Ideal first-gate candidate likelihood]
\label{prop:likelihood}
Fix an accepted public transcript of challenges and responses.  In the ideal
first-gate model, the likelihood of a candidate $s$ is zero if any support
inequality is violated.  Among feasible candidates it is proportional to
$P_s(G)^{-N}$, where $N$ is the number of signatures.
\end{proposition}

\begin{proof}
Before conditioning on $G$, a particular challenge has candidate-independent
probability and every mask vector has probability $(2Y+1)^{-768}$.  Thus a
particular feasible pair $(c,z)$ has the same unconditional mass for every
candidate; an infeasible pair has mass zero.  Conditioning each independent
record on $G$ divides its feasible mass by $P_s(G)$.  Multiplication over
$N$ records proves the claim.
\end{proof}

Proposition~\ref{prop:likelihood} says that maximum likelihood selects a
feasible candidate with the \emph{smallest} first-gate acceptance probability.
Using Lemma~\ref{lem:support}, its leading relative loss is
\begin{equation}
 \frac{1}{1{,}048{,}473}
 \E_c\sum_{r,j}\bigl(|(cs_r)_j|-51\bigr)_+.
\label{eq:hingeloss}
\end{equation}
For a random sparse signed challenge, a convolution coefficient is
approximately centered Gaussian with variance proportional to the squared
coefficient norm of the corresponding secret polynomial.  The expectation
of the hinge in Equation~\eqref{eq:hingeloss} increases with that variance.
This motivates maximizing squared norm over the feasible support polytope.

The norm replacement is a heuristic.  Proposition~\ref{prop:likelihood} is
exact only for the isolated ideal first gate, while genuine signatures have
later rejection gates and public hints.  We therefore calibrate the surrogate
on an exact first-gate simulator and separately test it on genuine signatures.
The experiments, rather than a theorem, establish recovery for the two keys
reported here.

\subsection{A public starting direction}

The retained bounds provide a useful first direction without any secret
labels.  For a normalized row, let $a=\langle x,s\rangle$.  Conditional on
that endpoint producing a row, the bound $b$ is uniform on the integers
\[
                       -495\le b\le\min(-51,a).
\]
Consequently,
\begin{equation}
              \E[2b+495\mid x,\text{row}]=\min(-51,a).
\label{eq:boundmoment}
\end{equation}
We regress the entirely public target $2b+495$ on $x$ with a fixed ridge of
$10^{-3}$ and retain only the normalized direction.  Although the censoring
in Equation~\eqref{eq:boundmoment} prevents this regression from being a
direct coefficient estimator, it gives a consistently useful radial start.

\section{Public recovery algorithm}
\label{sec:algorithm}

The final execution reads only the public transcript and public key.  Its
continuous stage is fixed for all six secret polynomials.

\begin{enumerate}[leftmargin=1.7em]
  \item Decode $z$ and $c$.  Convert every nonvacuous instance of
        Equation~\eqref{eq:rawconstraint} to
        $\langle x,s\rangle\ge b$, discarding the rest.
  \item Solve the public ridge regression from
        Equation~\eqref{eq:boundmoment} and normalize its result to a
        direction $d_0$.
  \item Initialize a working set with the 2,000 rows having the largest
        public bounds.  Solve
        \[
             \max_x\; d_i^T x
             \quad\text{subject to}\quad
             \langle x_r,x\rangle\ge b_r,\qquad -16\le x_j\le16.
        \]
        Check the result against every public row and add up to the 2,000
        largest violations.  Repeat separation until all continuous
        violations are at most $10^{-6}$.
  \item Set $d_{i+1}=x/\norm{x}_2$ and repeat the radial linear program at
        most ten times.  The working set is retained between iterations.
  \item Restrict every integer coefficient to
        $\{\lfloor x_j\rfloor,\lceil x_j\rceil\}\cap[-16,16]$.  For each
        polynomial, select the 5,000 rows with least public slack and run a
        binary integer program maximizing the frozen radial objective.  Give
        each polynomial 20 seconds, validate its incumbent against every
        exact support row, and add violations as cutting planes.  The six
        response-feasible incumbents form a public baseline $\bar s$.
  \item Recover $t_0$ as described in Section~\ref{sec:t0} and form
        $t=2^{10}t_1+t_0$.  Run the public singleton-support lattice ladder
        below for blocks $r=0,1,\ldots,5$ in that order.  Accept a candidate
        only after every collected support inequality and all 1,152 exact LWR
        residual bounds have been checked with integer arithmetic.  If the
        response-only baseline already passes both checks, it is itself a
        complete public recovery; the ladder may still be run as prescribed.
\end{enumerate}

The $10^{-6}$ tolerance is used only for floating-point separation.  Integer
support validation and the LWR relation use exact integer arithmetic.  In an
early genuine checkpoint, a $10^{-7}$ recheck repeatedly classified an active
LP boundary as violated at solver precision while adding no new row; the
larger tolerance and an explicit no-new-row check remove that numerical
stall without using secret information.

\subsection{Deterministic LWR block ladder}
\label{sec:lattice}

Fix a block $r$.  The other five polynomials remain equal to $\bar s$.  In
block $r$, put every coefficient at the lower endpoint of its public
floor/ceiling box and let $d\in\{0,1\}^v$ select the upper endpoint at the $v$
non-singleton coordinates.  Select $m=64$ LWR equations by a fixed public
order: enumerate seven-bit coefficient indices $i=0,1,\ldots,127$, replace
$i$ by its bit reversal, cycle through the nine output polynomials, and take
the first 64 rows.  After subtracting the fixed contribution, these equations
have the form
\begin{equation}
                         B d-b=e\pmod q,
             \qquad e\in[-16,15]^m.
\label{eq:blocklwr}
\end{equation}
Set $u=2d-\mathbf 1\in\{-1,1\}^v$ and
$h=2b-B\mathbf 1$.  Equation~\eqref{eq:blocklwr} implies
\begin{equation}
                         Bu-h=2e\pmod{2q}.
\label{eq:signedlwr}
\end{equation}
We use the row lattice generated by
\begin{equation}
 \mathcal B=
 \begin{pmatrix}
   2qI_m & 0             & 0\\
   B^T   & \lambda I_v   & 0\\
   h^T   & 0             & M
 \end{pmatrix},
 \qquad \lambda=18,\quad M=1.
\label{eq:embedding}
\end{equation}
If this block contains the needed corrections, the lattice contains
\begin{equation}
       (2e,\lambda u,-M),
\label{eq:planted}
\end{equation}
obtained by combining the secret rows with coefficients $u$, subtracting the
last row, and adding suitable $2qI_m$ rows.  For the successful block, the
embedding dimension is 187 and the vector in Equation~\eqref{eq:planted} has
squared norm 70,329.

Lattice reduction does not force the middle coordinates to be binary.  We run
LLL with $\delta=0.99$, followed when necessary by at most two BKZ-10 tours and
at most two BKZ-20 tours using a 256-bit MPFR Gram--Schmidt representation.  We inspect a
reduced row only if its final coordinate is exactly $\pm M$ and every middle
coordinate is exactly $\pm\lambda$.  After normalizing its sign, those middle
coordinates define $u$ and hence $d=(u+\mathbf1)/2$.  All other reduced rows
are ignored.  A decoded candidate still has no standing until it passes the
complete support and LWR validations.

Both reported runs fix the ladder to the six singleton supports.  On the
calibration key, blocks 0, 1, and 2 completed LLL, BKZ-10, and BKZ-20 without
a binary embedding row; block 3 succeeded after BKZ-10.  On the prospective
key, the integer baseline already passed every public check.  The prescribed
ladder was nevertheless run unchanged and rediscovered that identical
candidate for block 0 after BKZ-10.  A larger deterministic ladder could
continue to block pairs if all singletons failed; that extension was neither
needed nor tested here.

The optional post-experiment comparison with the generated secret is absent
from every regression, objective, row choice, coefficient box, lattice,
stopping rule, and public validation above.

\section{Recovering the omitted low part \texorpdfstring{$t_0$}{t0}}
\label{sec:t0}

The public key contains $t_1$ but Equation~\eqref{eq:lwr} needs the full
rounded product $t$.  The ATLAS hint supplies bounded public observations of
$t_0$.  From a signature, the verifier can compute
\begin{equation}
                 a=\operatorname{Round}(Az)-c(2^d t_1)\pmod p
\label{eq:hinta}
\end{equation}
and use the public hint to reconstruct the intended high part $w_1$.  Centering
the corresponding low representative gives an observable $u$ satisfying
\begin{equation}
                 u=ct_0+r_0,\qquad |r_0|\le8183.
\label{eq:hintobs}
\end{equation}
For each of the nine polynomials, we solve the public minimax problem
\begin{equation}
   \min_{x\in[-511,512]^{128}}\;\max_i |(c_i x-u_i)|
\label{eq:minimax}
\end{equation}
by a cutting-plane LP and round the solution.  A nearby-integer completion is
accepted only after every coefficient of every observation satisfies
Equation~\eqref{eq:hintobs}.

On each target used below, 20,000 signatures recovered all 1,152
coefficients.  On the calibration and prospective keys, respectively, the
continuous RMSE values were $0.116057$ and $0.120930$, and the maximum errors
were $0.438380$ and $0.464272$.  Both rounded vectors had maximum public
residual 8,183 and were returned unchanged by the independent integer pass.
Prospective collection took 123.9 seconds and the minimax solve took 260.9
seconds wall time with 482 MiB peak RSS.

This high-level partial-key recovery is analogous to the recovery of
Dilithium's omitted $t_0$ by Azevedo Oliveira, Calle Viera, Cogliati, and
Goubin \cite{Uncompressing}.  We do not claim the $t_0$ step alone as a new
attack.  Its role here is to expose the full public LWR validation relation
once the separate first-gate channel recovers $s_1$.

\section{Calibration and genuine-transcript experiments}
\label{sec:experiment}

\subsection{Ideal first-gate calibration}

Our ideal simulator samples a uniform signed weight-31 challenge, weights it
by its exact first-gate acceptance probability under the calibration secret,
and then samples every response coefficient uniformly from
$I((cs_1)_j)$.  It deliberately omits every later gate.  The public recovery
algorithm does not receive the calibration secret.

\begin{table}[ht]
\centering
\caption{Ideal first-gate calibration on the target secret.  ``Neighbors''
counts valid one-coordinate $\pm1$ alternatives that satisfy every support
row.}
\begin{tabular}{@{}rrrrr@{}}
\toprule
Signatures & Retained rows & Neighbors & Radial RMSE & Rounded exact\\
\midrule
100,000   & 64,327    & 807/1,497 & 12.4346  & 25/768\\
1,000,000 & 646,782   & 18/1,497  & 0.877129 & 344/768\\
3,000,000 & 1,939,568 & 0/1,497   & 0.238428 & 734/768\\
\bottomrule
\end{tabular}
\label{tab:ideal}
\end{table}

At three million ideal records, the maximum continuous coefficient error was
$0.857099$, placing every true integer inside its public floor/ceiling box.
The six 20-second integer programs recovered all 768 coefficients.  The
assembled result passed all 1,152 public LWR checks with residual range
$[-16,15]$.  Simulation took 511.8 seconds, radial recovery took 182.2 seconds
with 1.36 GiB peak RSS, and integer completion took 127.9 seconds with 1.48
GiB peak RSS.

\subsection{Experimental protocol}

Both genuine experiments used a four-core Intel Core i5-6500 desktop with
16~GiB RAM.  Distinct deterministic 64-byte key-generation seeds fix the two
ATLAS-128 keys for reproducibility.  For each key, query messages are distinct
16-byte encodings of the counters $0,\ldots,2{,}999{,}999$ under fixed public
byte pads.  Four processes handle disjoint consecutive counter ranges.  Every
process independently derives the same target key, calls the faithful signer,
parses the ordinary signature, and streams only the nonvacuous public rows.

For each key, the 20,000-record hint experiment uses the first 20,000 query
messages.  Because ATLAS signing is deterministic, these are repeated records
of queries already counted in the three-million set, rather than additional
distinct messages.  The public keys embedded by the hint and response
collectors match byte for byte.

The first experiment served as calibration during method development.  For
the second experiment, the complete procedure, parameters, block order, and
stopping rules were frozen before the next unused deterministic key identifier
was selected.  The freeze manifest has SHA-256 digest
\begin{center}
{\ttfamily\scriptsize \FreezeManifestHash}
\end{center}
The signing collectors necessarily held the target secret internally to
generate oracle responses, but their optional calibration outputs were
disabled and the recovery pipeline received only public artifacts.  Public
recovery, complete public validation, equivalent-key construction, and
fresh-message forgery all finished before the target secret was exported,
inspected, or supplied to any recovery or evaluation step.  A manifest of
those immutable public artifacts was then sealed with digest
\begin{center}
{\ttfamily\scriptsize \PublicOutcomeHash}
\end{center}
Only afterward did a separate calibration step regenerate the deterministic
secret and score the outputs.

\subsection{Calibration checkpoints and final recovery}

At 100,000 signatures, both laws are underdetermined.  At one million, the
genuine signer tracks the ideal calibration closely:

\begin{table}[ht]
\centering
\caption{Public radial estimator on ideal and genuine transcripts.}
\begin{tabular}{@{}lrrrr@{}}
\toprule
Transcript & Signatures & RMSE & Correlation & Rounded exact\\
\midrule
Ideal first gate & 100,000   & 12.4346  & 0.3312   & 25/768\\
Genuine signer   & 100,000   & 11.6469  & 0.4089   & 34/768\\
Ideal first gate & 1,000,000 & 0.877129 & --       & 344/768\\
Genuine signer   & 1,000,000 & 0.803564 & 0.996706 & 361/768\\
\bottomrule
\end{tabular}
\label{tab:checkpoints}
\end{table}

The genuine one-million transcript contains 644,828 rows.  Its public solve
took 93.4 seconds while four further collection processes ran and peaked at
518 MiB RSS.  Afterward, a calibration-only audit found 29 of 1,497 nearest
alternatives still feasible, compared with 18 in the ideal model.

The final genuine transcript contains \RealRows{} rows.  Its public radial
point has RMSE \RealRadialRMSE{} only in the later calibration comparison;
the recovery run did not know that score.  Direct rounding gives
\RealRoundedExact{} correct coefficients.  That comparison also shows that the
true integer lies in all 768 public floor/ceiling boxes.  Six response-only integer programs produce a
baseline satisfying every support row.  This baseline fails the public LWR
check, with none of 1,152 residuals in $[-16,15]$.  Only afterward, a
calibration comparison finds 745 correct coefficients: blocks 0, 1, 2, 4,
and 5 are exact, while block 3 has 23 one-step errors.

The immutable lattice-ladder output passes every support row and all 1,152
LWR checks.  A separate calibration comparison then confirms 768 of 768
coefficients, with maximum error zero.  The public singleton supports 0, 1,
and 2 consume their complete reduction schedules; support 3 succeeds after
BKZ-10.  The ladder takes \RealLatticeSeconds{} seconds and peaks at 590 MiB
RSS.  Its output, serialized as six lines of 128 signed decimal integers with
single spaces and newline terminators, has SHA-256 digest
\begin{center}
{\ttfamily\small \RecoveryHash}
\end{center}

This execution itself did not load the target secret, but the algorithm had
been developed while that secret was available.  In particular, the
single-block completion was designed after a diagnostic localized the
baseline's 23 errors to block 3.

\subsection{Prospectively frozen second-key recovery}

The second transcript contains \ProspectiveRows{} support rows and has
SHA-256 digest
\begin{center}
{\ttfamily\scriptsize \ProspectiveTranscriptHash}
\end{center}
The public radial point was written before calibration.  The later comparison
measured RMSE \ProspectiveRadialRMSE{}, maximum error
\ProspectiveRadialMaxError{}, correlation $0.999676153$, and
\ProspectiveRoundedExact{}/768 correct coefficients after direct rounding.
Every true coefficient lay in its public floor/ceiling box.

All six response-only integer incumbents passed every support inequality.
Their assembled baseline also passed all 1,152 public LWR checks with residual
range $[-16,15]$.  The prescribed lattice ladder was nevertheless run without
change; block 0 succeeded after BKZ-10 and returned the identical vector.
Before calibration, this vector had already formed an equivalent key and
produced an accepted fresh-message signature.  The later comparison found
exact $s_1$ recovery, 768/768, and exact $t_0$ recovery, 1,152/1,152.  A
post-experiment audit found that none of 1,481 valid one-coordinate $\pm1$
alternatives satisfied all support rows; every alternative violated at least
two.  Counting the separate 20,000-record hint recollection and every public
recovery stage, the measured prospective run took about 5,817 seconds
(97.0 minutes); a unified transcript collector could avoid the 123.9-second
hint recollection.  The serialized recovered $s_1$ has digest
\begin{center}
{\ttfamily\scriptsize \ProspectiveRecoveryHash}
\end{center}

\begin{table}[ht]
\centering
\small
\caption{Genuine-transcript recovery on the calibration and prospectively
frozen keys.  Calibration metrics in the lower rows were computed only after
the corresponding public artifacts were immutable.}
\begin{tabular}{@{}lrr@{}}
\toprule
Quantity & Calibration key & Prospective key\\
\midrule
Chosen-message signatures & 3,000,000 & 3,000,000\\
Retained support rows & \RealRows & \ProspectiveRows\\
Collection wall time (s) & \RealCollectionSeconds & \ProspectiveCollectionSeconds\\
Aggregate signatures/s & \RealCollectionRate & \ProspectiveCollectionRate\\
Continuous radial RMSE & \RealRadialRMSE & \ProspectiveRadialRMSE\\
Continuous maximum error & \RealRadialMaxError & \ProspectiveRadialMaxError\\
Correct after direct rounding & \RealRoundedExact/768 & \ProspectiveRoundedExact/768\\
Nearest alternatives surviving & 0/1,497 & 0/1,481\\
Radial solve wall time (s) & \RealRadialSeconds & \ProspectiveRadialSeconds\\
Integer baseline wall time (s) & \RealIntegerSeconds & \ProspectiveIntegerSeconds\\
Baseline exact coefficients & 745/768 & 768/768\\
Baseline public LWR check & failed & passed\\
Ladder stopping point & block 3, BKZ-10 & block 0, BKZ-10\\
Ladder wall time (s) & \RealLatticeSeconds & \ProspectiveLatticeSeconds\\
Peak recovery RSS & \RealPeakRSS & \ProspectivePeakRSS\\
Final exact $s_1$ coefficients & \RealRecoveredExact/768 & 768/768\\
Public LWR residual range & \RealLWRRange & $[-16,15]$\\
\bottomrule
\end{tabular}
\label{tab:final}
\end{table}
\FloatBarrier

\section{Equivalent signing keys and fresh-message forgeries}
\label{sec:forgery}

Recovering $s_1$ and $t_0$ does not reveal the original signing-randomness
seed $K$.  That seed is not needed for an equivalent key.  From the public
key we parse the public matrix seed $\rho$ and compute
$\mathit{tr}=H(\pk)$.  We choose a fresh $K'$, then pack
\begin{equation}
                  \sk'=(\rho,K',\mathit{tr},s_1,t_0)
\label{eq:equivkey}
\end{equation}
in the specified 2,128-byte secret-key format.

The forgery message is an 80-byte fixed string selected before the prospective
run.  Every signing-oracle query has length 16, so freshness follows
immediately, independently of message contents.  For each key, the faithful
signer using $\sk'$ returns a 2,081-byte detached signature, and the faithful
verifier using the corresponding original public key accepts it.  For
reproducibility, the SHA-256 digests are
\begin{center}
\begin{tabular}{@{}ll@{}}
Calibration signature: & {\ttfamily\scriptsize \ForgeryHash}\\
Calibration equivalent key: & {\ttfamily\scriptsize \EquivalentKeyHash}\\
Prospective signature: & {\ttfamily\scriptsize \ProspectiveForgeryHash}\\
Prospective equivalent key: & {\ttfamily\scriptsize \ProspectiveEquivalentKeyHash}
\end{tabular}
\end{center}

This establishes exact $s_1$ recovery followed by equivalent signing-key
construction and a successful fresh-message forgery for each tested key.  It
does not recover either original seed $K$.

\section{Scope, limitations, and interpretation}
\label{sec:scope}

The interval calculation in Lemma~\ref{lem:support} and every support row in
Equation~\eqref{eq:normalized} are exact consequences of the written
ATLAS-128 parameters.  The constant-width rejection argument in the
specification is therefore false whenever $|(cs_1)_j|>51$.  This is a design
issue, independently of the submitted repeated-mask implementation flaw.

The efficient decoder is partly heuristic.  Proposition~\ref{prop:likelihood}
is exact for the isolated ideal first gate, but the squared-norm surrogate,
Gaussian convolution model, and behavior under later signing gates are not
proved.  The two genuine experiments show that they work on both reported
targets.
They do not establish a key-averaged success probability or a three-million
worst-case bound.

The procedure was developed retrospectively on the calibration key.  A
secret-based diagnostic localized its 23 baseline errors to block 3 and
motivated the singleton ladder.  The final calibration execution used no
secret input and accepted only public checks, but its development could have
overfit that target.  For the second experiment, we froze and hashed the
complete procedure before generating a distinct deterministic key.  No
ground-truth secret file existed during recovery or forgery, the public result
was hashed before calibration, and the frozen procedure succeeded unchanged.

A formal EUF-CMA advantage is a probability over fresh key generation and the
adversary's randomness, which two successful experiments cannot estimate.
The concrete result is exact recovery and a fresh-message forgery on both
tested written-design keys, with one execution prospectively frozen.

The experiments cover ATLAS-128 only.  Other parameter sets have different
$n$, $\ell$, $\kappa$, $\eta$, and $\beta_1$, so neither the sample count nor
the same decoder should be extrapolated without a separate analysis.

\section{Conclusion}

ATLAS-128's first rejection gate intersects a translated response interval
with a finite mask support.  Because $\beta_1$ is much smaller than the
allowed challenge-product magnitude, that intersection is secret dependent.
The resulting support half-spaces do not identify the key by feasibility, but
their conditional likelihood favors large-radius feasible secrets.  On both
reported keys, public radial LPs narrow every coefficient to adjacent
integers, public integer completion recovers a response-feasible baseline,
and the complete LWR relation identifies exact $s_1$ from three million
genuine signatures.  The first key needs the public Kannan-embedding ladder;
the prospective second key is already exact after integer completion, and the
prescribed ladder independently rediscovers it.  Public hints recover $t_0$,
and equivalent keys sign fresh messages accepted by the original verifiers.

The result warrants independent human review and external reproduction.  It
also shows that rejection-sampling interval counts must include the original
sampler boundary: translation invariance of an unconstrained interval is not
enough.

\begin{thebibliography}{9}
\small

\bibitem{ATLAS}
S.~Adhikary, D.~Das, A.~De, A.~Ganguly, H.~Hart, A.~Karmakar,
S.~Kundu, C.~Li, P.~Mahadev, P.~Mondal, Q.~Norga, D.~Pal, and
K.~Srivatsan.
\newblock \emph{MORNING-ATLAS: A Fast and Lightweight Digital Signature from
Module Learning with Rounding}.
\newblock Algorithm Specifications and Supporting Documentation, version 1.0,
30 June 2026.
\newblock \url{https://cdn.jsdelivr.net/gh/ngcc-dev/ngcc-harness@main/sign-15/sign-15-spec.pdf}.

\bibitem{NGCC15}
ngcc.dev.
\newblock \emph{MORNING-ATLAS (sign-15) vulnerability report}.
\newblock Includes sign-15-4, ``Repeated masking coefficients expose the
signing key,'' credited to Xianhui Lu and Yijian Liu, 23 September 2026.
\newblock \url{https://ngcc.dev/reports/sign-15.html}.

\bibitem{Uncompressing}
P.~Azevedo Oliveira, A.~Calle Viera, B.~Cogliati, and L.~Goubin.
\newblock Uncompressing Dilithium's public key.
\newblock \emph{Cryptology ePrint Archive}, Paper 2024/1373, 2024;
published at CRYPTO 2025.
\newblock \url{https://eprint.iacr.org/2024/1373}.

\end{thebibliography}

\end{document}
